\begin{table}[t]
  \centering
  \scriptsize
  \caption{Transparent seed0 screening used to decide which additional baselines were promoted to matched three-seed evaluation. The early external-baseline screen predated the non-patent temporal probe; the later panel-expansion screen added \texttt{RF} and \texttt{GraphDTA-style} under the final five-model design.}
  \label{tab:recent-screening}
  \setlength{\tabcolsep}{3.0pt}
  \resizebox{\linewidth}{!}{
  \begin{tabular}{@{}lccccccp{5.3cm}@{}}
    \toprule
    Model & \texttt{patent\_temporal} & \texttt{nonpatent\_temporal} & \texttt{unseen\_drug} & \texttt{blind\_start} & \texttt{unseen\_target} & Promotion & Rationale \\
    \midrule
    \texttt{RF} &
    0.8249 / 0.7793 &
    0.7216 / 0.7648 &
    0.7279 / 0.8912 &
    0.4488 / 0.7589 &
    0.5618 / 0.7981 &
    main 3-seed promotion &
    strongest seed0 system on every screened temporal and synthetic axis except \texttt{unseen\_target}; therefore promoted across all main axes as the learner-control baseline \\
    \texttt{GraphDTA-style} &
    0.7741 / 0.7242 &
    0.6673 / 0.7444 &
    0.5386 / 0.8159 &
    0.4245 / 0.7426 &
    0.5289 / 0.7776 &
    targeted 3-seed promotion &
    architecture-diverse graph-family probe that is competitive enough on the main synthetic and temporal axes to test whether the expanded-panel conclusion depends on fixed-vector models \\
    \texttt{DTI-LM} &
    0.7897 / 0.7435 &
    -- &
    -- &
    0.3873 / 0.7130 &
    0.5308 / 0.8011 &
    main 3-seed promotion &
    strongest seed0 recent baseline on the primary real-OOD benchmark; therefore promoted in the main text \\
    \texttt{HyperPCM} &
    0.7686 / 0.7247 &
    -- &
    -- &
    0.3896 / 0.7447 &
    0.5294 / 0.7990 &
    secondary 3-seed promotion &
    competitive on synthetic seed0 axes, but weaker on the main patent benchmark; promoted only to test whether the main conclusion changes \\
    \bottomrule
  \end{tabular}}
\end{table}
